\section{Implémentation}

Cette section présente les détails techniques de l'implémentation de DesignPro AI, en se concentrant sur les aspects critiques du code et les décisions d'architecture.

\subsection{Structure du Code}

\subsubsection{Organisation des Fichiers}

\begin{lstlisting}[caption=Structure du projet]
mon-app-design/
├── apps/
│   └── web/
│       ├── app/                    # Next.js App Router
│       │   ├── (auth)/            # Routes authentification
│       │   ├── dashboard/         # Interface principale
│       │   └── api/               # API routes
│       ├── components/            # Composants React
│       │   ├── ui/               # Composants UI réutilisables
│       │   └── design/           # Composants spécifiques design
│       ├── lib/                   # Logique métier
│       │   ├── ai.ts             # Service IA (core)
│       │   ├── supabase.ts       # Client Supabase
│       │   └── utils.ts          # Utilitaires
│       └── styles/               # Styles globaux
├── packages/                      # Packages partagés (monorepo)
└── supabase/                     # Configuration Supabase
    ├── migrations/               # Migrations SQL
    └── seed.sql                  # Données de test
\end{lstlisting}

\subsubsection{Fichier Core : \texttt{lib/ai.ts}}

Le fichier \texttt{ai.ts} contient ~1000 lignes de code TypeScript structurées comme suit :

\begin{enumerate}[leftmargin=*]
    \item \textbf{Imports et Configuration} (lignes 1-50)
    \item \textbf{Types et Interfaces} (lignes 51-150)
    \item \textbf{Classe AIService} (lignes 151-1000)
    \begin{itemize}
        \item Phase 1 : Génération de prompts (lignes 200-350)
        \item Phase 2 : Génération d'images (lignes 351-600)
        \item Phase 3 : Agent Q (lignes 601-750)
        \item Phase 4 : R.E.A.L. (lignes 751-900)
        \item Utilitaires et fallbacks (lignes 901-1000)
    \end{itemize}
\end{enumerate}

\subsection{Implémentation des Phases}

\subsubsection{Phase 1 : Génération de Prompts}

\textbf{Fonction principale} :

\begin{lstlisting}[language=TypeScript, caption=Génération de prompt avec Mistral]
async generateProfessionalPrompt(
  projectData: any,
  methodology: string,
  methodologyParams?: Record<string, any>
): Promise<string> {
  try {
    // Construction du prompt système
    const systemPrompt = this.buildMethodologyPrompt(
      methodology, 
      projectData
    );
    
    // Appel Mistral AI
    const response = await fetch('https://api.mistral.ai/v1/chat/completions', {
      method: 'POST',
      headers: {
        'Authorization': `Bearer ${MISTRAL_API_KEY}`,
        'Content-Type': 'application/json'
      },
      body: JSON.stringify({
        model: 'mistral-large-latest',
        messages: [
          { role: 'system', content: systemPrompt },
          { role: 'user', content: `Générez un prompt pour: ${projectData.description}` }
        ],
        temperature: 0.7,
        max_tokens: 500
      })
    });
    
    const data = await response.json();
    return data.choices[0].message.content;
    
  } catch (error) {
    console.error('Erreur génération prompt:', error);
    return this.generateFallbackPrompt(projectData, methodology);
  }
}
\end{lstlisting}

\textbf{Gestion des erreurs} :

\begin{itemize}[leftmargin=*]
    \item Retry automatique (3 tentatives)
    \item Fallback vers templates pré-définis
    \item Logging détaillé pour debugging
\end{itemize}

\subsubsection{Phase 2 : Génération d'Images}

\textbf{Fonction principale} :

\begin{lstlisting}[language=TypeScript, caption=Génération d'images avec Stable Diffusion]
async generateDesignCandidates(
  prompt: string,
  modelId: string,
  settings: GenerationSettings
): Promise<DesignGenerationResult> {
  
  const models = [
    'stabilityai/stable-diffusion-xl-base-1.0',
    'stabilityai/stable-diffusion-2-1',
    'runwayml/stable-diffusion-v1-5'
  ];
  
  for (const model of models) {
    try {
      const result = await this.generateWithHuggingFace(
        model, 
        prompt, 
        settings
      );
      
      return {
        images: result.images,
        source: 'huggingface',
        model: model,
        used_fallback: false
      };
      
    } catch (error) {
      console.warn(`Échec modèle ${model}, tentative suivante...`);
      continue;
    }
  }
  
  // Fallback final : images de démonstration
  return {
    images: await this.generateRealisticFallbackImages(prompt),
    source: 'fallback',
    used_fallback: true,
    fallback_reason: 'Tous les modèles ont échoué'
  };
}
\end{lstlisting}

\textbf{Gestion du Cold Start} :

\begin{lstlisting}[language=TypeScript, caption=Gestion du cold start Hugging Face]
async generateWithHuggingFace(
  modelId: string, 
  prompt: string, 
  settings: any
) {
  let attempts = 0;
  const MAX_RETRIES = 3;
  
  while (attempts < MAX_RETRIES) {
    try {
      const response = await fetch(
        `https://router.huggingface.co/hf-inference/models/${modelId}`,
        {
          method: 'POST',
          headers: {
            'Authorization': `Bearer ${HF_API_TOKEN}`,
            'Content-Type': 'application/json'
          },
          body: JSON.stringify({
            inputs: prompt,
            parameters: {
              num_inference_steps: settings.quality,
              guidance_scale: settings.guidanceScale,
              width: settings.width,
              height: settings.height
            },
            options: { wait_for_model: true }
          })
        }
      );
      
      if (response.ok) {
        const blob = await response.blob();
        return { images: [await this.blobToBase64(blob)] };
      }
      
      // Gestion du cold start
      if (response.status === 503) {
        const errorData = await response.json();
        const waitTime = errorData.estimated_time || 20;
        
        console.log(`Modèle en chargement, attente ${waitTime}s...`);
        await new Promise(resolve => setTimeout(resolve, (waitTime + 2) * 1000));
        attempts++;
        continue;
      }
      
      throw new Error(`HTTP ${response.status}`);
      
    } catch (error) {
      attempts++;
      if (attempts >= MAX_RETRIES) throw error;
      await new Promise(resolve => setTimeout(resolve, 2000));
    }
  }
}
\end{lstlisting}

\subsubsection{Phase 3 : Agent Q}

\textbf{Analyse visuelle avec GPT-4 Vision} :

\begin{lstlisting}[language=TypeScript, caption=Analyse visuelle GPT-4 Vision]
private async analyzeDesignImageWithVision(
  imageUrl: string
): Promise<string> {
  
  if (OPENAI_API_KEY && OPENAI_API_KEY !== 'your-openai-api-key-here') {
    try {
      const response = await fetch(
        'https://api.openai.com/v1/chat/completions',
        {
          method: 'POST',
          headers: {
            'Authorization': `Bearer ${OPENAI_API_KEY}`,
            'Content-Type': 'application/json'
          },
          body: JSON.stringify({
            model: 'gpt-4-vision-preview',
            messages: [{
              role: 'user',
              content: [
                {
                  type: 'text',
                  text: this.buildVisionAnalysisPrompt()
                },
                {
                  type: 'image_url',
                  image_url: { 
                    url: imageUrl,
                    detail: 'high'
                  }
                }
              ]
            }],
            max_tokens: 500,
            temperature: 0.3
          })
        }
      );
      
      const data = await response.json();
      return data.choices[0].message.content;
      
    } catch (error) {
      console.error('Erreur GPT-4 Vision:', error);
      // Fallback vers Replicate ou description contextuelle
      return this.generateContextualDescription(imageUrl);
    }
  }
  
  return this.generateContextualDescription(imageUrl);
}
\end{lstlisting}

\subsection{Interface Utilisateur}

\subsubsection{Composants React Principaux}

\textbf{DesignWorkflow} : Composant principal orchestrant le workflow

\begin{lstlisting}[language=TypeScript, caption=Composant DesignWorkflow]
'use client';

import { useState } from 'react';
import { AIService } from '@/lib/ai';

export function DesignWorkflow() {
  const [phase, setPhase] = useState<1 | 2 | 3 | 4>(1);
  const [prompt, setPrompt] = useState('');
  const [images, setImages] = useState<string[]>([]);
  const [evaluation, setEvaluation] = useState<any>(null);
  
  const aiService = AIService.getInstance();
  
  const handlePhase1 = async () => {
    const generatedPrompt = await aiService.generateProfessionalPrompt(
      projectData,
      methodology
    );
    setPrompt(generatedPrompt);
    setPhase(2);
  };
  
  const handlePhase2 = async () => {
    const result = await aiService.generateDesignCandidates(
      prompt,
      selectedModel,
      settings
    );
    setImages(result.images);
    setPhase(3);
  };
  
  const handlePhase3 = async (imageUrl: string) => {
    const eval = await aiService.evaluateDesignWithAgentQ(
      imageUrl,
      prompt,
      methodology
    );
    setEvaluation(eval);
    setPhase(4);
  };
  
  return (
    <div className="workflow-container">
      {phase === 1 && <Phase1Component onNext={handlePhase1} />}
      {phase === 2 && <Phase2Component prompt={prompt} onNext={handlePhase2} />}
      {phase === 3 && <Phase3Component images={images} onNext={handlePhase3} />}
      {phase === 4 && <Phase4Component evaluation={evaluation} />}
    </div>
  );
}
\end{lstlisting}

\subsubsection{Styling avec TailwindCSS}

\begin{lstlisting}[language=HTML, caption=Exemple de composant stylé]
<div className="bg-gradient-to-br from-blue-50 to-indigo-100 rounded-xl p-6 shadow-lg">
  <h2 className="text-2xl font-bold text-gray-800 mb-4">
    Phase 1 : Génération de Prompt
  </h2>
  
  <div className="space-y-4">
    <select className="w-full px-4 py-2 border border-gray-300 rounded-lg focus:ring-2 focus:ring-blue-500">
      <option>Meuble</option>
      <option>Électronique</option>
      <option>Automobile</option>
    </select>
    
    <button className="w-full bg-blue-600 hover:bg-blue-700 text-white font-semibold py-3 px-6 rounded-lg transition-colors">
      Générer le Prompt
    </button>
  </div>
</div>
\end{lstlisting}

\subsection{Persistance des Données}

\subsubsection{Sauvegarde dans Supabase}

\begin{lstlisting}[language=TypeScript, caption=Sauvegarde de l'état du projet]
async saveProjectState(state: ProjectState) {
  const { data, error } = await supabase
    .from('project_states')
    .insert({
      user_id: user.id,
      project_id: state.projectId,
      methodology: state.methodology,
      prompt_generated: state.prompt,
      images: state.images,
      agent_q_evaluation: state.evaluation,
      real_simulation: state.simulation
    });
    
  if (error) {
    console.error('Erreur sauvegarde:', error);
    throw error;
  }
  
  return data;
}
\end{lstlisting}

\subsubsection{Chargement de l'Historique}

\begin{lstlisting}[language=TypeScript, caption=Chargement de l'historique]
async loadProjectHistory(projectId: string) {
  const { data, error } = await supabase
    .from('project_states')
    .select('*')
    .eq('project_id', projectId)
    .order('created_at', { ascending: false });
    
  if (error) throw error;
  return data;
}
\end{lstlisting}

\subsection{Optimisations}

\subsubsection{Caching}

\begin{itemize}[leftmargin=*]
    \item \textbf{Modèles Hugging Face} : Mis en cache côté serveur après premier chargement
    \item \textbf{Prompts générés} : Stockés en base pour réutilisation
    \item \textbf{Images} : Stockées dans Supabase Storage avec CDN
\end{itemize}

\subsubsection{Lazy Loading}

\begin{lstlisting}[language=TypeScript, caption=Lazy loading des composants]
import dynamic from 'next/dynamic';

const Phase3Component = dynamic(() => import('./Phase3'), {
  loading: () => <LoadingSpinner />,
  ssr: false
});
\end{lstlisting}

\subsubsection{Mode Low Memory}

Pour les utilisateurs avec GPU limité :

\begin{lstlisting}[language=TypeScript, caption=Mode low memory]
if (settings.lowMemoryMode) {
  pipeline = StableDiffusionPipeline.from_pretrained(
    modelId,
    torch_dtype=torch.float16,
    use_safetensors=True,
    variant="fp16"
  );
  pipeline.enable_attention_slicing();
  pipeline.enable_vae_slicing();
}
\end{lstlisting}

\subsection{Tests}

\subsubsection{Tests Unitaires}

\begin{lstlisting}[language=TypeScript, caption=Exemple de test]
import { AIService } from '@/lib/ai';

describe('AIService', () => {
  let aiService: AIService;
  
  beforeEach(() => {
    aiService = AIService.getInstance();
  });
  
  test('generateProfessionalPrompt retourne un prompt valide', async () => {
    const prompt = await aiService.generateProfessionalPrompt(
      { category: 'meuble', focus: 'ergonomie' },
      'DFX'
    );
    
    expect(prompt).toBeTruthy();
    expect(prompt.length).toBeGreaterThan(50);
    expect(prompt).toContain('ergonomic');
  });
  
  test('generateDesignCandidates retourne 4 images', async () => {
    const result = await aiService.generateDesignCandidates(
      'a modern chair',
      'stable-diffusion-xl',
      { numImages: 4 }
    );
    
    expect(result.images).toHaveLength(4);
    expect(result.source).toBe('huggingface');
  });
});
\end{lstlisting}

\subsubsection{Tests d'Intégration}

Tests du workflow complet end-to-end :

\begin{lstlisting}[language=TypeScript, caption=Test E2E]
test('Workflow complet Phase 1-4', async () => {
  // Phase 1
  const prompt = await aiService.generateProfessionalPrompt(...);
  expect(prompt).toBeTruthy();
  
  // Phase 2
  const images = await aiService.generateDesignCandidates(prompt, ...);
  expect(images.images.length).toBe(4);
  
  // Phase 3
  const evaluation = await aiService.evaluateDesignWithAgentQ(
    images.images[0], 
    prompt, 
    'DFX'
  );
  expect(evaluation.overall_score).toBeGreaterThan(0);
  expect(evaluation.overall_score).toBeLessThanOrEqual(10);
  
  // Phase 4
  const simulation = await aiService.simulateREALAnalysis(...);
  expect(simulation.fea_analysis).toBeDefined();
  expect(simulation.dfm_analysis).toBeDefined();
});
\end{lstlisting}

\subsection{Déploiement}

\subsubsection{Configuration Vercel}

\begin{lstlisting}[language=JSON, caption=vercel.json]
{
  "framework": "nextjs",
  "buildCommand": "npm run build",
  "devCommand": "npm run dev",
  "installCommand": "npm install",
  "env": {
    "NEXT_PUBLIC_SUPABASE_URL": "@supabase-url",
    "NEXT_PUBLIC_SUPABASE_ANON_KEY": "@supabase-anon-key",
    "MISTRAL_API_KEY": "@mistral-api-key",
    "HF_API_TOKEN": "@hf-api-token",
    "OPENAI_API_KEY": "@openai-api-key"
  },
  "regions": ["iad1"],
  "functions": {
    "app/api/**/*.ts": {
      "maxDuration": 60
    }
  }
}
\end{lstlisting}

\subsubsection{CI/CD Pipeline}

\begin{lstlisting}[language=YAML, caption=.github/workflows/deploy.yml]
name: Deploy to Vercel

on:
  push:
    branches: [main]

jobs:
  deploy:
    runs-on: ubuntu-latest
    steps:
      - uses: actions/checkout@v2
      
      - name: Install dependencies
        run: npm install
      
      - name: Run tests
        run: npm test
      
      - name: Build
        run: npm run build
      
      - name: Deploy to Vercel
        uses: amondnet/vercel-action@v20
        with:
          vercel-token: ${{ secrets.VERCEL_TOKEN }}
          vercel-org-id: ${{ secrets.ORG_ID }}
          vercel-project-id: ${{ secrets.PROJECT_ID }}
\end{lstlisting}

Cette implémentation robuste garantit la fiabilité, la performance et la maintenabilité de DesignPro AI en production.
